\section{}

\textbf{Tiempo de lectura:} \{s11-reading-time\}

La \textbf{memoria compartida}{} es una estrategia para comunicar
procesos donde uno de ellos gana acceso a regiones de la memoria del
otro; algo que por lo general el sistema operativo siempre intenta
evitar. Por eso, para que pueda haber memoria compartida es necesario
que los dos procesos estén de acuerdo en eliminar dicha restricción.

Dos procesos que comparten una región de la memoria pueden intercambiar
información simplemente leyendo y escribiendo datos en la misma. Sin
embargo debemos tener en cuenta que:

\begin{itemize}
\item
  La estructura de los datos y su localización dentro de la región
  compartida la determinan los procesos en comunicación y no el sistema
  operativo, a diferencia de lo que ocurre en los sistemas de paso de
  mensajes.
\item
  Los procesos son responsables de sincronizarse para no escribir y leer
  en el mismo sitio de la memoria al mismo tiempo, pues esto puede
  generar inconsistencias (véase el \hyperref[sincronizaciuxf3n]{???}) .
\end{itemize}

Las principales ventajas de la memoria compartida frente a otros
mecanismos de comunicación son:

\begin{itemize}
\item
  \textbf{Eficiencia}. Puesto que la comunicación tiene lugar a la
  velocidad de la memoria principal, se trata de un mecanismo
  tremendamente rápido.
\item
  \textbf{Conveniencia}. Puesto que el mecanismo de comunicación solo
  requiere leer y escribir de la memoria, se trata de un sistema muy
  sencillo y fácil de utilizar.
\end{itemize}

Como ocurre con las tuberías (véase el \hyperref[_tuberuxedas]{???}) la
memoria compartida puede ser anónima o con nombre.

\section{Memoria compartida
anónima}\label{_memoria_compartida_anuxf3nima}

La \textbf{memoria compartida anónima}{} solo existe para el proceso que
la crea y para sus procesos hijos, que heredan el acceso. Es por tanto,
una forma eficiente de comunicar procesos padres e hijos.

En los sistemas POSIX, las funciones y operadores de reserva de memoria
como \{clang\_malloc\} y \{cpp\_new\}, utilizan internamente la llamada
al sistema \{linux\_mmap\}. Esta función se puede llamar de la siguiente
manera para reservar \texttt{length} bytes de memoria.

\begin{Shaded}
\begin{Highlighting}[]
\DataTypeTok{void}\OperatorTok{*}\NormalTok{ p }\OperatorTok{=}\NormalTok{ mmap}\OperatorTok{(} 
\NormalTok{    NULL}\OperatorTok{,}
\NormalTok{    length}\OperatorTok{,}     
\NormalTok{    PROT\_READ }\OperatorTok{|}\NormalTok{ PROT\_WRITE}\OperatorTok{,}      
\NormalTok{    MAP\_ANONYMOUS }\OperatorTok{|}\NormalTok{ MAP\_PRIVATE}\OperatorTok{,} 
    \OperatorTok{{-}}\DecValTok{1}\OperatorTok{,}
    \DecValTok{0}
\OperatorTok{);}
\end{Highlighting}
\end{Shaded}

\begin{itemize}
\item
  Si todo va bien, \{linux\_mmap\} devuelve un puntero al primer byte de
  la memoria reservada.
\item
  Cantidad de memoria a reservar en bytes.
\item
  Permisos de acceso para la memoria reservada. En este caso, se
  solicita permitir la lectura y la escritura de la memoria.
\item
  \texttt{MAP\_ANONYMOUS} indica que la memoria no está respaldada por
  ningún archivo, por lo que su contenido inicial será cero. Mientras
  que \texttt{MAP\_PRIVATE} establece que la región de memoria es
  privada.
\end{itemize}

Lo interesante es que si se cambia \texttt{MAP\_PRIVATE} por
\texttt{MAP\_SHARED} la región de memoria reservada es memoria
compartida:

\begin{Shaded}
\begin{Highlighting}[]
\DataTypeTok{void}\OperatorTok{*}\NormalTok{ p }\OperatorTok{=}\NormalTok{ mmap}\OperatorTok{(}
\NormalTok{    NULL}\OperatorTok{,}
\NormalTok{    length}\OperatorTok{,}
\NormalTok{    PROT\_READ }\OperatorTok{|}\NormalTok{ PROT\_WRITE}\OperatorTok{,}
\NormalTok{    MAP\_ANONYMOUS }\OperatorTok{|}\NormalTok{ MAP\_SHARED}\OperatorTok{,} 
    \OperatorTok{{-}}\DecValTok{1}\OperatorTok{,}
    \DecValTok{0}
\OperatorTok{);}
\end{Highlighting}
\end{Shaded}

\begin{itemize}
\item
  Memoria anónima y compartida.
\end{itemize}

Es decir, que al crear un hijo con \{linux\_fork\} este tendrá una copia
de toda la memoria del proceso padre, excepto esta región en particular,
que será la misma que la del padre. Por lo tanto, escribiendo y leyendo
en esa región, ambos procesos pueden comunicarse.

En \{anom\_shared\_memory\_cpp\} se puede ver un ejemplo muy simple,
similar a \{fork\_pipe\_cpp\} pero utilizando memoria compartida para
comunicar ambos procesos. Como se puede apreciar, la versión que usa
memoria compartida es bastante más sencilla que la que utiliza tuberías.

En Microsoft Windows se puede hacer algo similar con
\{win32\_createfilemapping\}:

\begin{Shaded}
\begin{Highlighting}[]
\NormalTok{HANDLE hMapFile }\OperatorTok{=}\NormalTok{ CreateFileMapping}\OperatorTok{(} 
\NormalTok{    INVALID\_HANDLE\_VALUE}\OperatorTok{,}
\NormalTok{    NULL}\OperatorTok{,}
\NormalTok{    PAGE\_READWRITE}\OperatorTok{,} 
    \DecValTok{0}\OperatorTok{,}
\NormalTok{    length}\OperatorTok{,}         
\NormalTok{    NULL}
\OperatorTok{);}
\end{Highlighting}
\end{Shaded}

\begin{itemize}
\item
  Permisos de acceso para la memoria reservada.
\item
  Cantidad de memoria a reservar en bytes.
\item
  A diferencia de \{linux\_mmap\}, \{win32\_createfilemapping\} crea un
  objeto de memoria compartida, pero no hace visible esa memoria para
  nuestro proceso. Para eso hay que llamar a \{win32\_mapviewoffile\}
  pasándole el manejador \texttt{hMapFile} devuelto por
  \{win32\_createfilemapping\}.
\end{itemize}

\section{Memoria compartida con
nombre}\label{_memoria_compartida_con_nombre}

La \textbf{memoria compartida con nombre}{} es pública para el resto del
sistema, por lo que teóricamente cualquier proceso con permisos puede
acceder a ella para comunicarse con otros procesos.

Como ocurre en las tuberías con nombre, los \textbf{objetos de memoria
compartida con nombre} hay que crearlos antes de comenzar a utilizarlos.
Para eso los sistemas POSIX ofrecen la función \{linux\_shm\_open\}.

\begin{Shaded}
\begin{Highlighting}[]
\DataTypeTok{int}\NormalTok{ shmfd }\OperatorTok{=}\NormalTok{ shm\_open}\OperatorTok{(}   
    \StringTok{"/foo{-}shm"}\OperatorTok{,}         
\NormalTok{    O\_RDWR }\OperatorTok{|}\NormalTok{ O\_CREAT}\OperatorTok{,}   
    \BaseNTok{0666}                
\OperatorTok{);}
\end{Highlighting}
\end{Shaded}

\begin{itemize}
\item
  Nombre que identifica al objeto de memoria compartida. Como ocurre con
  los archivos, varios procesos pueden acceder al mismo objeto indicando
  el mismo nombre.
\item
  Valores que indican diferentes opciones a la hora de abrir el objeto.
  Por ejemplo, usando \texttt{O\_RDWR} indicamos que se abra para
  lectura y escritura. Mientras que con \texttt{O\_CREAT} se indica que
  el objeto debe crearse si no existía previamente.
\item
  Indica los permisos del objeto de memoria compartida al crearlo nuevo,
  de forma similar a los permisos que se aplican a los archivos en el
  sistema de archivos.
\end{itemize}

El valor devuelto por \{linux\_shm\_open\} es el descriptor del objeto
de memoria compartida, que utilizaremos posteriormente con
\{linux\_mmap\} al reservar una región de la memoria de nuestro proceso
donde ese objeto de memoria compartida será visible:

\begin{Shaded}
\begin{Highlighting}[]
\DataTypeTok{void}\OperatorTok{*}\NormalTok{ p }\OperatorTok{=}\NormalTok{ mmap}\OperatorTok{(}
\NormalTok{    NULL}\OperatorTok{,}
\NormalTok{    length}\OperatorTok{,}                 
\NormalTok{    PROT\_READ }\OperatorTok{|}\NormalTok{ PROT\_WRITE}\OperatorTok{,}
\NormalTok{    MAP\_SHARED}\OperatorTok{,}             
\NormalTok{    shmfd}\OperatorTok{,}                  
    \DecValTok{0}                       
\OperatorTok{);}
\end{Highlighting}
\end{Shaded}

\begin{itemize}
\item
  Al pasar el descriptor del objeto de memoria compartida, ya no se
  puede indicar \texttt{MAP\_ANONYMOUS}.
\item
  Se puede hacer visible para el proceso todo el objeto de memoria
  compartida o solo una parte. Para esto último se indica el tamaño de
  la región y el desplazamiento dentro del objeto, que es el último
  argumento de \{linux\_mmap\}.
\end{itemize}

Un objeto de memoria compartida recién creado tiene tamaño 0. Para
redimensionarlo se utiliza \{linux\_ftruncate\}, que lo que necesita es
el descriptor del objeto y el nuevo tamaño.

En \{shared\_memory\_cpp\} se puede ver el ejemplo de un programa que
muestra periódicamente la hora del sistema. En este caso controlado por
\{shared\_memory\_control\_cpp\} mediante memoria compartida.

En Microsoft Windows también se utiliza \{win32\_createfilemapping\}
para crear el objeto de memoria compartida con nombre. Simplemente hay
que indicar el nombre en el último argumento de la función.

\begin{Shaded}
\begin{Highlighting}[]
\NormalTok{HANDLE hMapFile }\OperatorTok{=}\NormalTok{ CreateFileMapping}\OperatorTok{(}
\NormalTok{    INVALID\_HANDLE\_VALUE}\OperatorTok{,}
\NormalTok{    NULL}\OperatorTok{,}
\NormalTok{    PAGE\_READWRITE}\OperatorTok{,}
    \DecValTok{0}\OperatorTok{,}
\NormalTok{    length}\OperatorTok{,}
    \StringTok{"Global}\SpecialCharTok{\textbackslash{}\textbackslash{}}\StringTok{FooMemoriaCompartida"} 
\OperatorTok{);}
\end{Highlighting}
\end{Shaded}

\begin{itemize}
\item
  Nombre del nuevo objeto de memoria compartida.
\end{itemize}
